% New capabilities from KACTL; original source code is unchanged.
\section{Bellman--Ford: affected destinations}
This variant labels each destination affected by a negative cycle, instead of only returning a cycle flag.
\lstinputlisting[style=cpp]{kactl/content/graph/BellmanFord.h}
\section{Biconnected edge blocks}
This variant returns edge IDs via a callback. Bridges are not emitted by the original implementation.
\lstinputlisting[style=cpp]{kactl/content/graph/BiconnectedComponents.h}
\section{Recovering a minimum s-t cut}
The original source explains extraction from residual reachability after maximum flow; it does not define another flow algorithm.
\lstinputlisting[style=cpp]{kactl/content/graph/MinCut.h}
\section{Global minimum cut}
\lstinputlisting[style=cpp]{kactl/content/graph/GlobalMinCut.h}
\section{Gomory--Hu tree}
\lstinputlisting[style=cpp]{kactl/content/graph/GomoryHu.h}
\section{Matching in general graphs}
\lstinputlisting[style=cpp]{kactl/content/graph/GeneralMatching.h}
\section{Edge coloring}
\lstinputlisting[style=cpp]{kactl/content/graph/EdgeColoring.h}
\section{Enumerating maximal cliques}
\lstinputlisting[style=cpp]{kactl/content/graph/MaximalCliques.h}
\section{Maximum clique}
\lstinputlisting[style=cpp]{kactl/content/graph/MaximumClique.h}
\section{Maximum independent set}
The original source explains the reduction to MaximumClique; it does not define a separate function.
\lstinputlisting[style=cpp]{kactl/content/graph/MaximumIndependentSet.h}
\section{Dynamic forests: link-cut tree}
\lstinputlisting[style=cpp]{kactl/content/graph/LinkCutTree.h}
\section{Directed minimum spanning arborescence}
\lstinputlisting[style=cpp]{kactl/content/graph/DirectedMST.h}
